\chapter{Conclusion and Summary}
\label{chap:conclusion}

This thesis investigated whether decoupling representation learning from downstream regression improves molecular optical property prediction in data-scarce and distribution-shifted settings. Specifically, we evaluated hybrid pipelines in which a directed message passing neural network (D-MPNN) encoder is first trained end-to-end on the downstream task, then frozen and paired with a post-hoc gradient-boosted tree regressor. We compared this approach against end-to-end fine-tuning, training from scratch, and several classical and physics-informed baselines across multiple datasets and split strategies. Several broad patterns emerged.

First, in small-$N$ regimes for absorption prediction, training a post-hoc XGBoost model on learned representations yields measurable improvements over end-to-end training at negligible additional computational cost. This effect is particularly pronounced under scaffold-based distribution shift, where training and test molecules differ structurally. In these regimes, the explicit structural regularization of the tree-based model can mitigate the instability of the high-capacity neural decoder, which is otherwise prone to overfitting. This hybrid advantage was less systematic for emission prediction, suggesting that static ground-state graphs may lack the necessary information to reliably capture excited-state relaxation and solvent reorganization without end-to-end adaptation.

Second, variance decomposition analyses isolated the mechanisms driving this hybrid advantage. While encoder initialization and training dynamics account for the dominant portion of performance variance across all data regimes, the variability introduced by the downstream regression head becomes substantial in small-$N$ settings. By replacing the highly variable neural head with a regularized ensemble of trees, the hybrid architecture appears to learn a smoother mapping from the unstable learned representations to the response.

Third, transfer learning experiments clarified the boundaries of pretraining utility. Fine-tuning pretrained encoders provides substantial performance gains when the downstream task aligns chemically with the pretraining corpus. However, under severe distribution shift, pretraining offers negligible or even negative returns. In successful transfer scenarios, freezing the fine-tuned encoder and applying a hybrid readout preserves the adapted latent structure while maximizing downstream statistical efficiency.

Beyond predictive performance, analysis of the learned latent space demonstrates that the D-MPNN encoder captures relatively stable, chemically meaningful structure without explicit supervision. The dominant embedding dimension strongly correlates with effective $\pi$-conjugation length and the HOMO--LUMO gap, while secondary dimensions disentangle structural and environmental factors like molecular size and polarity.

These results suggest a practical modelling guideline. In experimental photophysics regimes where data are scarce or distribution shift is anticipated, decoupling representation learning from regression is a beneficial strategy. Freezing a task-trained graph encoder and fitting a gradient-boosted tree on the embeddings provides competitive or improved performance for virtually no additional overhead once the encoder is trained. 

Conversely, in large-$N$ settings, fully joint optimization becomes sufficiently constrained and the benefit of decoupling diminishes. Future work may extend this analysis to properties more strongly governed by excited-state relaxation, alternative encoder architectures, or domain adaptation techniques specifically designed for severe distribution shift. Nonetheless, the central finding remains clear: the physically grounded representations learned by graph neural networks can be effectively leveraged by simpler, regularized regressors, offering a practical modelling strategy for data-limited chemical domains.